\labsection{1}{Directories and process creation}{sec:lab01}
\begin{labproject}{Explore a directory and create a child process}
\textbf{Coverage.} Chapters 1--2.\\
\textbf{Source files.} \texttt{source\_code/Lab01/task\_1.c}, \texttt{task\_2.c}.\\
\textbf{Goal.} Observe two fundamental OS abstractions: the directory namespace and process identity.

\begin{labmode}
The supplied programs are intentionally small. Treat them as instrumentation: run them, inspect the output, then change one thing at a time so you can connect a C API call to an observable OS effect.
\end{labmode}

\medskip\noindent\textbf{Task A --- Directory traversal.}\par
\begin{lstlisting}[style=osc]
DIR *dirp = opendir(path);
struct dirent *dptr;
while ((dptr = readdir(dirp)) != NULL) {
    printf("%s\n", dptr->d_name);
}
closedir(dirp);
\end{lstlisting}

\textbf{Observe.} Entries such as \texttt{.} and \texttt{..} may appear. The listing order should not be assumed to be alphabetic unless your program sorts it explicitly.

\medskip\noindent\textbf{Task B --- Parent and child PIDs.}\par
\begin{lstlisting}[style=osc]
pid_t pid = fork();
if (pid == 0)
    printf("child: %ld\n", (long)getpid());
else if (pid > 0)
    printf("parent: %ld\n", (long)getpid());
\end{lstlisting}

\begin{tryit}
Run the process program several times. Does parent output always appear before child output? Add \texttt{wait(NULL)} in the parent and explain how the ordering changes.
\end{tryit}
\end{labproject}

\begin{labdeliverable}
Submit terminal output for both tasks plus a short explanation of (1) what \texttt{readdir()} returns, (2) the three possible \texttt{fork()} return cases, and (3) why output order can vary without \texttt{wait()}.
\end{labdeliverable}
